% ============================================================
% Section 9: Conclusion (Medium length for TNNLS)
% Target: ~1 page
% ============================================================

\section{Conclusion}
\label{sec:conclusion}

We studied when a controlled graph-aggregation operator improves or harms class-conditional discriminability, and how far simple empirical diagnostics transfer to trained models. The result is a scoped mechanism study with negative diagnostic evidence, not a universal model-selection theory.

In the binary fixed-degree conditional-neighbor model, the exact ratio $\kappa=\rho^2d/[1+(1-\rho^2)s]$ identifies how edge correlation, degree, and feature strength jointly change one-hop Mahalanobis discriminability. The center--neighbor expression covers a prescribed linear surrogate. Across a frozen grid of 2,000 configuration--seed records, empirical moments match the exact expression with 0.94\% median absolute relative error over nonzero ratios. This validates the implemented moment calculation within the surrogate, not a broadcasting threshold or trained-GNN accuracy.

The empirical diagnostics remain deliberately weaker. ADR is a matched percentage-point difference, and the one-hop/two-hop quantities describe archived label structure; neither identifies a winning architecture. The audited selector comparison shows no incremental value over an always-graph baseline, and $1-\mathrm{Acc}_{\mathrm{MLP}}$ remains only arithmetic positive-gain headroom. The prospective selector benchmark is not yet executed, while the legacy edge-randomization outcomes have not been regenerated with the verified degree-matched implementation.

\subsection{Theoretical Significance and Scope}

The central theoretical contribution is the scoped moment-based discriminability analysis, not a universal guarantee for trained GNNs. Finite graphs, learned transformations, dependent neighborhoods, and optimization can shift the predicted transition.

\subsection{Limitations and Future Work}

The exact propositions assume binary labels, fixed degree, homogeneous conditional neighbor-label probabilities, conditionally independent neighbors, and Gaussian features with shared covariance. They provide no scalar multi-class extension and no random-degree or finite-depth GNN guarantee.

The next empirical priorities are a clean prospective diagnostic benchmark and a degree-matched edge-randomization rerun with concurrent structural statistics. Theoretical extensions to random degree, dependent neighborhoods, attention, heterogeneous graphs, and temporal graphs remain separate work.

\subsection{Reproducibility}

\balance

An anonymized code, result, and supplementary-material repository is provided through the submission system. The public URL will be restored after anonymous review.

The audited release contains the frozen direct-validation configuration, 2,000 immutable records, deterministic summary, and figure. It also contains prospective diagnostic scoring and verified edge-randomization code, but no new model outcomes from those two pending reruns. Historical exploratory outputs are retained only with explicit provenance limitations and are not claimed as fully reproducible evidence.
